\documentclass[conference]{IEEEtran}
\IEEEoverridecommandlockouts
\usepackage{cite}
\usepackage{amsmath,amssymb,amsfonts}
\usepackage{algorithmic}
\usepackage{graphicx}
\usepackage{textcomp}
\usepackage{xcolor}
\usepackage{booktabs}
\def\BibTeX{{\rm B\kern-.05em{\sc i\kern-.025em b}\kern-.08em
    T\kern-.1667em\lower.7ex\hbox{E}\kern-.125emX}}

\begin{document}

\title{GraphAIOps: Unifying Troubleshooting Guide Generation,
Multi-Agent Triage, and Graph-Based Root Cause Analysis
for Cloud Incident Management}

\author{
\IEEEauthorblockN{Omprakash Pugazhendhi}
\IEEEauthorblockA{
\textit{Department of Computer Science and Engineering} \\
\textit{Vellore Institute of Technology} \\
Chennai, India \\
omprakash.2021@vitstudent.ac.in
}
}

\maketitle

\begin{abstract}
Managing incidents in large-scale cloud systems is one of the most operationally
demanding challenges in modern software engineering.
On-call engineers are paged at all hours, expected to diagnose and fix outages
quickly, and are often working from troubleshooting guides that are outdated,
incomplete, or simply missing.
This paper presents \textbf{GraphAIOps}, a unified synthesis of three recent
Microsoft Research systems that together automate the full incident lifecycle:
TSGen generates structured troubleshooting guides directly from historical
incident telemetry; TRIANGLE routes each incoming alert to the correct engineering
team through a multi-agent negotiation mechanism; and a Causal Knowledge Graph
with graphlet inference surfaces root causes from years of post-mortem records,
augmented by a fine-tuned language model.
We show that graph-based reasoning is the technical thread that connects all three
systems, and that treating them as a single integrated pipeline---rather than
independent tools---unlocks compounding improvements that none of them achieves alone.
\end{abstract}

\begin{IEEEkeywords}
cloud reliability, AIOps, incident management, knowledge graph, multi-agent LLM,
troubleshooting guide generation, root cause analysis, graphlet inference
\end{IEEEkeywords}


\section{Introduction}

Picture an on-call engineer paged at 3\,a.m. because a latency spike is affecting
checkout on a production service. The first thing they reach for is the
troubleshooting guide---but the guide was last updated eight months ago, references
a dashboard that no longer exists, and was written for a slightly different version
of the service. Meanwhile, the incident has already been routed to the wrong team
by an automated classifier that misread the error signature, and re-routing it
takes another twenty minutes. By the time the right engineer is looking at the right
logs, the outage has lasted over an hour.

This scenario is not hypothetical. Cloud services at hyperscale generate thousands
of incidents daily, and the three failure modes described above---stale documentation,
mis-triage, and slow root cause analysis---are the dominant reasons why mean time
to resolution (MTTR) remains stubbornly high even at mature engineering organizations.
Each failure mode has attracted its own line of research: AutoTSG~\cite{autotsq}
automated the execution of existing troubleshooting guides; DeCaf~\cite{decaf}
applied machine learning to KPI-based incident routing; RCACopilot~\cite{rcacopilot}
introduced chain-of-thought LLM reasoning for root cause analysis.
What none of these systems does is share artifacts or feed outputs from one stage
into the next.

This paper synthesizes three recent Microsoft Research systems---TSGen~\cite{tsgen},
TRIANGLE~\cite{triangle}, and KG-Graphlet RCA~\cite{ckg}---into a unified pipeline
we call \textbf{GraphAIOps}. Each system was designed independently to solve one
of the three failure modes above. We argue that when you look at them together,
the architectural choices converge: all three represent knowledge as graphs,
all three use graph traversal or graph attention to answer queries over that
knowledge, and all three produce artifacts that the other two systems can consume.
That convergence is not a coincidence---it reflects something real about the
structure of incident data, which is fundamentally relational and does not fit
naturally into flat retrieval or tabular classifiers.

The rest of this paper describes each system in turn (Sections~III--V),
then shows how they integrate into a self-reinforcing pipeline (Section~VI),
and closes with a discussion of what remains unsolved (Section~VII).


\section{Related Work}

The incident management literature is large, but prior work almost uniformly
addresses one stage at a time. On the documentation side, Nissist~\cite{nissist}
converts existing human-written TSGs into knowledge graphs that can answer
step-by-step queries; it improves TSG usability but does not solve the problem
of TSGs that are missing or outdated to begin with. StepFly~\cite{stepfly} ran
an empirical study on 92 real TSGs and found that \textit{instruction drift}---
where an LLM agent ignores the guide structure and rewrites queries on the fly---
is the primary failure mode in agentic TSG execution. LLexus~\cite{llexus}
generates executable automation plans from human-authored TSGs, moving closer to
full autonomy but still requiring humans to write the source material. FLASH
addresses the same execution problem using a ReAct-based LLM agent with dynamic
context injection during runtime.

On the triage and diagnosis side, DeCaf~\cite{decaf} uses machine learning and
pattern mining over KPI time-series data to triage incidents, without any natural
language understanding of incident tickets. RCACopilot~\cite{rcacopilot} goes
further by applying GPT-4 with chain-of-thought prompting to generate root cause
hypotheses from alert signals, but it does not exploit the relational structure
of historical post-mortem data. To our knowledge, no existing system unifies
automated guide generation, multi-team triage with shared knowledge, and
graph-structured root cause analysis into a single integrated pipeline.


\section{TSGen: Synthesizing Troubleshooting Guides from Incident History}

An internal Microsoft study examined over 4,000 troubleshooting guides and found
a clear pattern: when a high-quality, up-to-date TSG exists and an engineer finds
it, incident mitigation time drops by roughly 40\%~\cite{tsgen}. The catch is
that most TSGs do not meet that bar. They are written once, rarely maintained,
stored in silos that engineers don't know to look in, and often cover the happy
path without addressing the edge cases that cause the most damage in production.

TSGen attacks this problem at the source by generating TSGs from the raw material
that accumulates naturally as incidents are resolved: IcM tickets, diagnostic log
snippets, Kusto query results, and the resolution notes left by engineers who
fixed the problem. The system processes this material through five stages.
\textit{Collection} gathers the relevant signals for a given monitor or service
using monitor IDs or custom Kusto queries.
\textit{Filtering} applies machine learning to remove noise---duplicate signals,
false positives, entries from unrelated incidents---and surface the attributes
most predictive of resolution path.
\textit{Core Incident Selection} picks a small, representative subset of incidents
that together cover the distinct failure modes associated with a given service,
avoiding redundancy while ensuring coverage.
\textit{Data Distillation} reads the resolution paths of those core incidents and
extracts the diagnostic checks and mitigation steps that actually worked.
\textit{TSG Generation} assembles the distilled steps into a structured,
action-oriented guide using a large language model, formatted so that both human
engineers and automated triage agents can consume it directly.

That last point matters for the integrated pipeline. Because TSGen's output is
machine-readable and consistently structured, TRIANGLE agents can be trained on
the generated corpus without any additional formatting work. In pilot deployments
at Microsoft, dozens of AI-generated TSGs passed review and were published for
production on-call use---a high bar, since engineers are under no obligation to
accept machine-generated documentation and tend to be skeptical of it.


\section{TRIANGLE: Multi-Agent Negotiation for Incident Triage}

Getting an incident to the wrong team is expensive. The wrong team spends time
reading a ticket that isn't theirs, figures out it belongs elsewhere, and
re-routes it. Time-to-engage---the interval between an alert firing and the
right engineer starting to work---is the metric that takes the hit, and in a
serious outage that delay is measured in user impact.

TRIANGLE approaches triage as a multi-agent reasoning problem~\cite{triangle}.
Each engineering team has an agent that encodes that team's domain knowledge
through two sources: its historical incident corpus and the TSGs generated by
TSGen. When a new incident arrives, the system does not try to classify it
with a single model trained on a static snapshot of team boundaries. Instead,
it asks each agent whether the incident belongs to its team, collects the
reasoning alongside the vote, and runs a negotiation phase in which agents
revise their positions in light of each other's arguments until consensus emerges.

Three problems make naive classification fail here. \textit{Semantic heterogeneity}
means that the same root cause---say, a connection pool exhaustion in a shared
database---can look completely different in text depending on which surface it
manifests on and who filed the ticket. \textit{Decentralized knowledge} means
that the correct owner of an incident is often not determinable from any single
team's perspective alone; a network misconfiguration that causes latency on a
storage service is technically a networking team issue even if the storage team
sees it first. \textit{Dynamic responsibilities} means that team ownership changes
as organizations restructure, and any static classifier trained three months ago
may already be wrong. The negotiation mechanism handles all three: by letting
agents reason over shared artifacts and revise their conclusions, the system
adapts to ambiguity and organizational change without retraining.

In production, TRIANGLE achieves up to 97\% triage accuracy and has reduced
time-to-engage by up to 91\% across a platform serving tens of millions of users.


\section{KG-Graphlet: Graph-Structured Root Cause Analysis}

Once the right team is engaged, they still face the hardest part: figuring out
why the incident happened. The answer is almost always somewhere in the
organization's accumulated history of post-mortem documents---Problem Review
Board (PRB) records that describe past incidents, their root causes, and how they
were resolved. The problem is that this knowledge is buried in natural language,
spread across thousands of documents, and essentially unsearchable with keyword
queries for anything beyond the most common failure modes~\cite{ckg}.

The approach taken here builds a \textit{Causal Knowledge Graph} (CKG) over the
full PRB corpus by extracting entities---symptoms, root causes, resolutions, and
services involved---and connecting them with typed edges that encode causal
and temporal relationships. This turns an unstructured document archive into a
queryable graph where semantic clustering over node embeddings reveals structure:
in practice, 60 distinct symptom clusters emerge, each grouping incidents that
share underlying failure patterns even when the surface descriptions are very different.

The most technically interesting contribution in this line of work is
\textit{graphlet inference}. Rather than finding nodes similar to the current
incident, the system searches for recurring small subgraph patterns---graphlets---
that capture how failures propagate through the service dependency graph.
A graphlet might encode the pattern: ``service A's error rate spikes, database B's
connection count hits its ceiling 90 seconds later, and downstream service C
begins returning 503s.'' This pattern can match even if A, B, and C have different
names in the new incident, as long as the structural relationship is the same.
Graphlet matching generalizes from experience in a way that node-level similarity
cannot, because it captures the \textit{mechanism} of failure rather than its
surface appearance.

A GPT-3.5 model fine-tuned on prompts constructed from CKG subgraphs improves
root cause generation quality by 45.5\% over zero-shot prompting, and mitigation
suggestion quality by 131.3\%. In a live production evaluation, over 70\% of
on-call engineers rated the system's recommendations 3 out of 5 or higher---
meaningful, given that experienced engineers are deeply skeptical of automated RCA
and set a high bar for what counts as useful.


\section{GraphAIOps: Integration and Feedback}

\begin{figure}[t]
\centering
\renewcommand{\arraystretch}{1.4}
\begin{tabular}{c}
\hline
\textbf{IcM Alert Fired} \\
\hline
\textbf{Stage 1 --- TRIANGLE (Triage)} \\
semantic distillation $\rightarrow$ agent negotiation \\
$\Rightarrow$ route to responsible team \\[4pt]
\textbf{Stage 2 --- KG-Graphlet (Root Cause)} \\
CKG query $\rightarrow$ graphlet matching \\
$\Rightarrow$ root cause hypothesis surfaced \\[4pt]
\textbf{Stage 3 --- TSGen (Resolution)} \\
corpus lookup or on-demand TSG synthesis \\
$\Rightarrow$ structured guide retrieved or generated \\[4pt]
\textbf{Engineer Confirms and Executes Mitigation} \\
\hline
\textit{Resolved incident $\rightarrow$ new PRB record $\rightarrow$} \\
\textit{feeds back into CKG, TSGen, and TRIANGLE} \\
\hline
\end{tabular}
\caption{The GraphAIOps end-to-end pipeline. Each resolved incident
produces artifacts that improve all three stages in the next cycle.}
\label{fig:pipeline}
\end{figure}

Figure~\ref{fig:pipeline} shows the full pipeline. When an alert fires, TRIANGLE
routes it to the right team in seconds rather than hours. The CKG is immediately
queried using the incident's entity signatures, and the graphlet matcher surfaces
a root cause hypothesis along with the historical incidents that support it.
TSGen's corpus is searched for the most relevant guide; if a close match exists,
it is retrieved; if not, TSGen synthesizes one from similar resolved incidents
on demand. The engineer reviews the hypothesis and the guide, then executes or
adjusts the suggested mitigation.

What makes this arrangement more than a pipeline is the feedback loop that closes
after each resolution. A resolved incident produces a new PRB record, which the
CKG ingests to enrich its graphlet library. The resolution steps feed back into
TSGen as new training material, keeping the guide corpus fresh without any manual
effort. The updated TSGs flow back into the knowledge base that TRIANGLE agents
draw on for the next triage cycle. The system does not require separate maintenance
teams for each component---it improves through use.

Table~\ref{tab:impact} reports the measured impact of each component from
independent production evaluations. The numbers do not compound trivially, but
even independently they represent substantial improvements to incident response.

\begin{table}[h]
\caption{Measured Production Impact of Each Component}
\label{tab:impact}
\centering
\begin{tabular}{lp{3.2cm}r}
\toprule
\textbf{System} & \textbf{Metric} & \textbf{Result} \\
\midrule
TSGen       & Mitigation time reduction  & $-40\%$ \\
TRIANGLE    & Triage accuracy            & up to $97\%$ \\
TRIANGLE    & Time-to-engage reduction   & $-91\%$ \\
KG-Graphlet & Root cause quality gain    & $+45.5\%$ \\
KG-Graphlet & Mitigation suggestion gain & $+131.3\%$ \\
KG-Graphlet & Engineer satisfaction ($\geq$3/5) & $>70\%$ \\
\bottomrule
\end{tabular}
\end{table}


\section{Discussion}

The shared architectural choice across all three systems---representing incident
knowledge as a graph rather than a flat corpus---is not arbitrary. Incident data
is fundamentally relational. An alert does not have meaning in isolation; it has
meaning in the context of which services depend on which other services, which
teams own which components, and which failure patterns have occurred before under
similar conditions. Flat retrieval treats all of that context as metadata to be
discarded or encoded as feature vectors. Graph representations preserve it in a
form that supports structured reasoning.

Several open problems remain. \textit{Cold-start} is genuine: a new service with
no incident history has no PRB records in the CKG, no resolved incidents for TSGen
to learn from, and no base for TRIANGLE's agent. Transfer learning across services
and organizations could address this, but the right inductive bias for
cross-service failure pattern transfer is an open research question.
\textit{Agent drift} in TRIANGLE is also real: as team responsibilities change,
an agent's domain knowledge becomes stale, and retraining requires curating a fresh
corpus of accepted and rejected incidents---a non-trivial operational burden.
Finally, both the CKG and TSGen are evaluated on proprietary Microsoft corpora,
which makes independent benchmarking impossible for the research community.
Publishing an anonymized incident dataset with synthetic entity labels would be a
significant contribution to the field, independent of any system paper.

Looking ahead, the most impactful direction is federated incident learning: sharing
graphlet patterns across organizations without exposing raw telemetry. A smaller
cloud operator running the same database software as a hyperscaler could benefit
from failure propagation patterns discovered at scale, even if the specific service
names differ. The graph representation used by the CKG is well-suited to this kind
of abstracted knowledge transfer in a way that raw log sharing is not.


\section{Conclusion}

GraphAIOps brings together three production AI systems that each solve a distinct
but interconnected problem in cloud incident management. TSGen ensures that
troubleshooting guides exist, are current, and are structured for both humans and
machines. TRIANGLE ensures that the right team engages with an incident immediately,
without the round-trips that mis-triage causes. KG-Graphlet inference gives that
team a concrete root cause hypothesis grounded in the organization's full history
of past incidents, not just the engineer's personal experience.

What this paper argues is that these three systems are stronger together than
apart---not just because their outputs compose cleanly, but because each system's
artifacts improve the others in a cycle. Resolved incidents become better TSGs,
better CKG entries, and better agent training data. The pipeline learns from every
incident it handles, and over time it should close the gap between the median
engineer's response and the best engineer's response. That is a genuinely ambitious
goal for operational AI, and the production results from Microsoft suggest it is
within reach.

\begin{thebibliography}{00}

\bibitem{tsgen}
C.~Bansal et al.,
``TSGen: Automated Troubleshooting Guide Generation from Cloud Incidents,''
in \textit{Proc. FSE}, Montreal, Canada, Jul. 2026.

\bibitem{triangle}
Z.~Yu, M.~Ma, X.~Feng, R.~Ding, C.~Zhang, Z.~Li, M.~Chintalapati, X.~Zhang,
R.~Wang, C.~Bansal, S.~Rajmohan, Q.~Lin, S.~Zhang, C.~Pei, and D.~Pei,
``TRIANGLE: Automated Incident Triaging for Cloud Reliability Data,''
in \textit{Proc. ASE}, 2025.

\bibitem{ckg}
C.~Bansal et al.,
``Mining Root Cause Knowledge from Cloud Service Incident Investigations for AIOps,''
\textit{arXiv:2204.11598}, 2022.

\bibitem{autotsq}
M.~Shetty, C.~Bansal, S.~P.~Upadhyayula, A.~Radhakrishna, and A.~Gupta,
``AutoTSG: Machine Learning and Rule Based Framework for Automated Troubleshooting,''
in \textit{Proc. ESEC/FSE}, 2022.

\bibitem{rcacopilot}
X.~Chen et al.,
``RCACopilot: On-Call Incident Root Cause Analysis with LLMs,''
in \textit{Proc. EuroSys}, 2024.

\bibitem{decaf}
M.~Ma et al.,
``DeCaf: Diagnosing and Triaging Performance Issues in Large-Scale Cloud Services,''
in \textit{Proc. ICSE}, 2020.

\bibitem{nissist}
Z.~Guo et al.,
``Nissist: An Incident Mitigation Copilot Based on Troubleshooting Guides,''
\textit{arXiv:2402.17531}, 2024.

\bibitem{stepfly}
X.~Liu et al.,
``StepFly: An Agentic Framework for Troubleshooting Guide Automation,''
\textit{arXiv:2510.10074}, 2025.

\bibitem{llexus}
``LLexus: LLM Agents for Executable TSG Plans,''
\textit{Tech. Report, Microsoft Research}, 2024.

\end{thebibliography}

\end{document}
